\documentclass[11pt,a4paper]{article}
\usepackage[utf8]{inputenc}
\usepackage[T1]{fontenc}
\usepackage[spanish]{babel}
\usepackage{geometry}
\geometry{a4paper, margin=2.2cm}
\usepackage{xcolor}
\usepackage{tikz}
\usetikzlibrary{shapes.geometric, arrows.meta}
\usepackage{array}
\usepackage{hyperref}
\usepackage{amsmath}
\usepackage{titlesec}
\usepackage{fancyhdr}
\usepackage{helvet}
\renewcommand{\familydefault}{\sfdefault}

% Color Palette: TutorMate Pro
\definecolor{primary}{HTML}{002B66}     % Dark Blue
\definecolor{accent}{HTML}{0D9488}      % Turquoise
\definecolor{gold}{HTML}{D97706}        % Gold
\definecolor{lightbg}{HTML}{F8FAFC}     % Light Gray

\hypersetup{
    colorlinks=true,
    linkcolor=accent,
    filecolor=accent,      
    urlcolor=accent,
    pdftitle={Guía de Inicio - TutorMate Pro},
    pdfpagemode=FullScreen,
}

% Page styling
\pagestyle{fancy}
\fancyhf{}
\rhead{\color{accent}\textbf{TutorMate Pro}}
\lhead{\color{primary}\small\leftmark}
\rfoot{\color{primary}\thepage}
\lfoot{\color{primary}\scriptsize tutormatepro.com --- Acompañamiento Individualizado}
\renewcommand{\headrulewidth}{0.4pt}
\renewcommand{\footrulewidth}{0.4pt}

% Section styling
\titleformat{\section}
  {\color{primary}\normalfont\Large\bfseries}
  {\thesection}{1em}{}
\titleformat{\subsection}
  {\color{accent}\normalfont\large\bfseries}
  {\thesubsection}{1em}{}

% TikZ Styles for Flowchart
\tikzset{
    block/.style={
        rectangle, 
        rounded corners=5pt, 
        minimum width=5cm, 
        minimum height=1.1cm, 
        text centered, 
        draw=primary, 
        fill=lightbg, 
        line width=1.2pt,
        text width=4.8cm
    },
    arrow/.style={
        thick,
        -Stealth,
        color=accent,
        line width=1.5pt
    }
}

\begin{document}

% --- PORTADA ---
\begin{titlepage}
    \centering
    \vspace*{2cm}
    {\Huge \color{primary} \textbf{TutorMate Pro}} \\[0.5cm]
    {\huge \color{accent} \textbf{Pack 4 clases}} \\[0.8cm]
    {\Large \color{gold} \textbf{Guía de Inicio y Bienvenida}} \\[1.5cm]
    
    \begin{minipage}{0.8\textwidth}
        \centering
        \large \color{primary} Acompañamiento individualizado para resolución rápida de dudas
    \end{minipage} \\[2cm]
    
    \vspace*{\fill}
    {\large \color{primary} \textbf{Autor: Nahuel García}} \\
    {\large \color{primary} tutormatepro.com} \\
    {\small \color{accent} \today}
\end{titlepage}

\newpage
\tableofcontents
\newpage

\section{Bienvenida al Programa}
¡Te damos la bienvenida a \textbf{TutorMate Pro}! Al adquirir este programa, has dado un paso fundamental hacia el dominio y la comprensión profunda de las matemáticas y la estadística. 

Este es un programa online de acompañamiento con \textbf{clases en vivo 1:1 (individuales)}. No se trata de un curso grabado estático: es un servicio diseñado a tu medida, donde dispondrás de un tutor dedicado que te guiará paso a paso, adaptándose a tus necesidades académicas específicas, tu ritmo y tus metas particulares.

Este \textbf{Pack de 4 clases} está diseñado para estudiantes que necesitan un apoyo puntual y concentrado. Es la opción perfecta para resolver dudas complejas en vísperas de un examen, recibir diagnóstico de cara a un tema específico o resolver atascos operativos rápidos en un temario definido.

\section{¿Qué incluye este Plan?}
Este paquete de acompañamiento individualizado te proporciona las siguientes herramientas y recursos clave:
\begin{itemize}
    \item 4 sesiones online individuales en vivo.
    \item Explicación y desglose de ejercicios paso a paso.
    \item Acceso a pizarra digital interactiva compartida.
    \item Orientación y diagnóstico rápido de bases.
    \item Apuntes teóricos detallados en PDF tras cada sesión.
\end{itemize}

\section{¿Para quién es este Programa?}
Este programa está estructurado para adaptarse con la máxima flexibilidad a estudiantes en distintas etapas académicas:
\begin{itemize}
    \item \textbf{Educación Secundaria (ESO):} Afianzamiento de bases operativas, superación del bloqueo con el álgebra básica y fomento de la seguridad en la resolución de problemas.
    \item \textbf{Bachillerato y Selectividad (EBAU/PAU):} Preparación rigurosa de exámenes trimestrales, asimilación del análisis matemático (cálculo), la probabilidad y la geometría espacial.
    \item \textbf{Universidad:} Acompañamiento en asignaturas de alto nivel como Álgebra Lineal, Cálculo Infinitesimal y Métodos Numéricos.
    \item \textbf{Estadística Aplicada:} Comprensión teórica y desarrollo práctico de estadística descriptiva e inferencial para diversas carreras universitarias.
    \item \textbf{Preparación de Convocatorias:} Refuerzo de verano, preparación de exámenes extraordinarios o nivelación pre-curso.
\end{itemize}

\newpage
\section{Cómo se organiza el Trabajo}
Para garantizar que cada minuto de clase sea productivo, el programa sigue una metodología organizada y personalizada. Antes del inicio o durante las primeras interacciones, el tutor recopila la siguiente información del estudiante:
\begin{enumerate}
    \item \textbf{Curso y Nivel Académico:} Para enmarcar el rigor matemático adecuado.
    \item \textbf{Temario o Guía Docente:} El temario específico oficial de su institución.
    \textbf{Fecha de Examen:} Para planificar el ritmo y priorizar temas clave.
    \item \textbf{Ejercicios y Exámenes Anteriores:} El material práctico que el alumno debe dominar.
    \item \textbf{Dificultades Principales:} Conceptos específicos donde se encuentra bloqueado.
\end{enumerate}
Con base en esta información, se define una \textbf{ruta de trabajo personalizada} para optimizar el número de sesiones contratadas.

\section{Ejemplo de Ruta de Trabajo}
La siguiente tabla muestra una planificación genérica orientativa del paquete. Es importante destacar que, al ser clases individuales, esta ruta se adapta y puede variar sustancialmente para ajustarse a tu ritmo real y prioridades académicas.

\begin{table}[h]
\centering
\renewcommand{\arraystretch}{1.3}
\begin{tabular}{|l|p{11cm}|}
\hline
\rowcolor{primary}
\color{white}\textbf{Sesión} & \color{white}\textbf{Enfoque de Trabajo} \\ \hline
 Clase 1 & Diagnóstico inicial de dificultades, priorización de dudas del temario y resolución de los primeros ejercicios clave. \\ \hline
 Clase 2 & Desarrollo teórico y práctico guiado de los temas con mayor peso en el examen. \\ \hline
 Clase 3 & Resolución de problemas de nivel avanzado y simulación de tipos de examen. \\ \hline
 Clase 4 & Sesión de cierre, corrección de fallos típicos, consejos de tiempo y dudas finales. \\ \hline
\end{tabular}
\caption{Planificación orientativa de las sesiones}
\end{table}

\newpage
\section{Metodología TutorMate Pro}
Nuestra filosofía didáctica se aleja de la memorización mecánica de fórmulas y procedimientos. Nos enfocamos en construir un aprendizaje sólido bajo las siguientes pautas:
\begin{itemize}
    \item \textbf{Entender antes de memorizar:} Explicar el "por qué" detrás de cada teorema, fórmula o algoritmo matemático.
    \item \textbf{Detección y corrección de errores de base:} La mayoría de los problemas en niveles avanzados no se deben al tema actual, sino a fallos acumulados en álgebra básica o aritmética. Nos detenemos a resolver estas lagunas.
    \item \textbf{Resolución estructurada paso a paso:} Enseñamos al alumno a desglosar problemas complejos en pasos sencillos, mejorando su redacción y organización matemática (clave para obtener el máximo puntaje en exámenes de selectividad y universidad).
    \item \textbf{Conexión entre teoría y práctica:} Trabajamos directamente sobre ejemplos concretos e ilustrativos para afianzar la teoría abstracta.
\end{itemize}

\begin{figure}[h]
\centering
\begin{tikzpicture}[node distance=1.6cm, auto]
  \node (val) [block] {\textbf{1. Valoración y Diagnóstico} \\ \scriptsize Definición de metas y nivel inicial};
  \node (plan) [block, below of=val] {\textbf{2. Ruta Personalizada} \\ \scriptsize Planificación según temario y fechas};
  \node (class) [block, below of=plan] {\textbf{3. Clases Guiadas 1:1} \\ \scriptsize Pizarra digital interactiva y explicación};
  \node (prac) [block, below of=class] {\textbf{4. Práctica de Consolidación} \\ \scriptsize Boletines de problemas adaptados};
  \node (rev) [block, below of=prac] {\textbf{5. Cierre y Simulación} \\ \scriptsize Consejos de examen y revisión final};
  
  \draw [arrow] (val) -- (plan);
  \draw [arrow] (plan) -- (class);
  \draw [arrow] (class) -- (prac);
  \draw [arrow] (prac) -- (rev);
\end{tikzpicture}
\caption{Esquema metodológico de aprendizaje TutorMate Pro}
\label{fig:flujo}
\end{figure}

\newpage
\section{Recomendaciones para el Alumno}
Para maximizar el aprovechamiento de tus clases individuales, te recomendamos adoptar las siguientes prácticas:
\begin{itemize}
    \item \textbf{Envía el material con antelación:} Comparte tu guía docente, apuntes o boletines de ejercicios antes de la clase para que el tutor los tenga analizados.
    \item \textbf{Apunta tus dudas concretas:} Lleva una lista con los puntos exactos donde te has atascado en tu estudio personal.
    \item \textbf{Prepara tu espacio de estudio:} Dispón de papel, bolígrafo, calculadora y una conexión a internet estable.
    \item \textbf{Trabaja entre sesiones:} Las clases online son el espacio para guiarte y aclarar dudas, pero la asimilación real ocurre cuando resuelves los ejercicios de forma autónoma entre clases.
    \item \textbf{Planifica con tiempo:} No acumules todas las clases para los dos días previos al examen; espaciar las sesiones permite asimilar los conceptos mucho mejor.
\end{itemize}

\section{Temarios Disponibles para Trabajar}
A continuación, se listan las principales materias y áreas que se pueden abordar en el acompañamiento, clasificadas por niveles académicos:

\subsection{Educación Secundaria (ESO)}
Operaciones fraccionarias, ecuaciones de primer y segundo grado, sistemas de ecuaciones lineales, funciones básicas, proporcionalidad directa e inversa, áreas y volúmenes geométricos, estadística descriptiva elemental.

\subsection{Bachillerato y EBAU / PAU}
\begin{itemize}
    \item \textbf{Análisis:} Límites de funciones, continuidad, cálculo de derivadas, optimización, integrales inmediatas y por métodos de integración.
    \textbf{Álgebra Lineal:} Matrices, operaciones, cálculo de determinantes y resolución de sistemas lineales (método de Rouché-Frobenius y regla de Cramer).
    \item \textbf{Geometría en el Espacio ($R^3$):} Posiciones relativas de rectas y planos, cálculo de distancias, ángulos y áreas en el espacio.
    \item \textbf{Probabilidad y Estadística:} Probabilidad condicionada, Teorema de Bayes, distribución Binomial y distribución Normal.
\end{itemize}

\subsection{Nivel Universitario y Estadística}
\begin{itemize}
    \item \textbf{Álgebra Avanzada:} Espacios vectoriales, bases, aplicaciones lineales, diagonalización de matrices y cálculo de bloques de Jordan.
    \item \textbf{Cálculo Infinitesimal:} Cálculo multivariable, integración múltiple, ecuaciones diferenciales de primer orden y sucesiones.
    \item \textbf{Estadística Descriptiva e Inferencial:} Intervalos de confianza, contrastes de hipótesis (paramétricos y no paramétricos), análisis de varianza (ANOVA), regresión lineal simple y múltiple, y su interpretación.
\end{itemize}

\section{Condiciones Generales del Programa}
El desarrollo del programa de acompañamiento se rige por las siguientes pautas profesionales:
\begin{itemize}
    \item \textbf{Formato de impartición:} Todas las sesiones son 100\% online en vivo, utilizando pizarra interactiva y videoconferencia.
    \item \textbf{Coordinación de horarios:} Las fechas de las clases se coordinarán directamente entre el tutor y el alumno según disponibilidad de agendas.
    \item \textbf{Política de reprogramación:} Si necesitas mover una clase, debes avisar al tutor con al menos 24 horas de antelación para poder reorganizar la sesión sin penalización.
    \item \textbf{Compromiso académico:} Las clases representan un soporte especializado y una guía didáctica de alto nivel, pero los resultados del alumno dependen directamente de su estudio independiente, constancia y realización de ejercicios prácticos recomendados.
\end{itemize}

\section{Próximo Paso}
Para comenzar con tu programa, por favor coordina tus horarios iniciales compartiendo tu temario, boletines de ejercicios actuales o la fecha de tus próximos exámenes. Con base en esta información estructuraremos tu ruta de trabajo a medida para aprovechar al máximo cada una de tus sesiones.

\end{document}
